%===================================================================
\subsection{State-Space Reduction via Size-Bin Moments}
\label{sec:size_bin_moments}
%===================================================================

This section presents the hybrid discrete + size-bin moment framework
used in RadCluster\_1\_0. The full per-size master equations
(Sec.~\ref{sec:ME_SIA}--\ref{sec:ME_He}) are replaced by a much
smaller ODE system in which the smallest cluster sizes are still
resolved one ODE per size (\emph{discrete range}), while the larger
sizes are aggregated into logarithmically spaced bins
(\emph{moment range}). Various grouping methods have been advanced
during the past decades
\cite{Kiritani1973,Koiwa1974,Hayns1976,Ghoniem1988BiasFactor,
HuangGhoniem1995,Golubov2001Grouping,Ovcharenko2003}. Three intra-bin
closures are supported in the code (\texttt{bin\_moment\_rates.py}):
piecewise-constant ($P=1$ moment per bin), hat-function / linear
dual basis ($P=2$), and log-normal ($P=3$). \textbf{Both the SIA and
the vacancy populations are reduced in the same way}; the He state is
attached to the vacancy bins (Case 1) or carried as a scalar total
(Case 2). The total ODE count drops from $O(I+V)$ to
\begin{equation}\label{eq:Neq_bin}
    N_{\mathrm{eq}}
    \;=\; \underbrace{i_{\mathrm{discrete}} + P\,I_{\mathrm{bin}}}_{\text{SIA block}}
    \;+\; \underbrace{v_{\mathrm{discrete}} + P\,V_{\mathrm{bin}}}_{\text{vacancy block}}
    \;+\; n_{\mathrm{He}}
    \;+\; 5,
\end{equation}
where the trailing $5$ counts the cumulative-loss accounting integrals
(Sec.~\ref{sec:mass_conservation}). Typically 20--60 bins suffice for
$I, V = 10^4\text{--}10^5$ with $r \in [1.5, 2]$.

%-------------------------------------------------------------------
\subsubsection*{Hybrid Bin Partition}
\label{sec:log_bins}
%-------------------------------------------------------------------

The SIA size axis $\{1, 2, \ldots, I\}$ is partitioned into a
\emph{discrete} prefix $\{1, \ldots, i_{\mathrm{discrete}}\}$ that is
tracked one ODE per size, followed by $I_{\mathrm{bin}}$
geometrically spaced bins that cover
$\{i_{\mathrm{discrete}}{+}1, \ldots, I\}$. The discrete cutoff
$i_{\mathrm{discrete}}$ is a numerical parameter chosen independently
of the physical mobility cutoff $n_{\max}^i$; in particular,
$i_{\mathrm{discrete}}$ may exceed $n_{\max}^i$ if extra resolution at
small immobile loops is desired. The bin edges
$\{s_k\}_{k=0}^{I_{\mathrm{bin}}}$ satisfy
\begin{equation}\label{eq:bin_boundaries}
    s_0 = i_{\mathrm{discrete}} + 1,
    \qquad
    s_{I_{\mathrm{bin}}} \;\le\; I + 1,
\end{equation}
and are generated by the recursion
\begin{equation}\label{eq:bin_ratio}
    s_{k+1} \;=\; \max\!\bigl(\,\lfloor r\,s_k\rfloor,\; s_k + 1\,\bigr),
    \qquad r > 1,
\end{equation}
which guarantees at least one integer per bin even when $r\,s_k$
floors back to $s_k$. The bin ratio is determined from the user's
target bin count $I_{\mathrm{bin}}$:
\begin{equation}\label{eq:r_from_Ibin}
    r \;=\; \max\!\Bigl(\,
      \bigl(I / i_{\mathrm{discrete}}\bigr)^{1/I_{\mathrm{bin}}},
      \; 1.01\,\Bigr).
\end{equation}
Bin $k$ spans the integer interval
$\mathcal{B}_k = \{n \in \mathbb{Z} : s_k \le n < s_{k+1}\}$ and has
width $\Delta_k = s_{k+1} - s_k$. The geometric midpoint of bin $k$ is
\begin{equation}\label{eq:shat}
    \hat{s}_k = \sqrt{s_k\,(s_{k+1}-1)}.
\end{equation}
The vacancy axis is partitioned identically with parameters
$v_{\mathrm{discrete}}$, $V_{\mathrm{bin}}$, edges
$\{t_k\}_{k=0}^{V_{\mathrm{bin}}}$, and ratio
$r_v = (V/v_{\mathrm{discrete}})^{1/V_{\mathrm{bin}}}$.

\paragraph{Example.} For $I = 5000$, $i_{\mathrm{discrete}} = 100$,
$I_{\mathrm{bin}} = 6$: $r \approx 2$ and 106 equations replace the
4\,900 per-size SIA ODEs.

%-------------------------------------------------------------------
\subsubsection*{Bin Moments}
\label{sec:bin_moments}
%-------------------------------------------------------------------

For each bin $k$ of the SIA partition, define
\begin{equation}\label{eq:bin_moments_def}
    \mu_k^{(0)}(t) \equiv \sum_{n \in \mathcal{B}_k} c_n(t),
    \qquad
    \mu_k^{(1)}(t) \equiv \sum_{n \in \mathcal{B}_k} n\,c_n(t),
    \qquad
    \mu_k^{(2)}(t) \equiv \sum_{n \in \mathcal{B}_k} n^2\,c_n(t).
\end{equation}
The mean cluster size in bin $k$ is
$\langle n\rangle_k = \mu_k^{(1)}/\mu_k^{(0)}$, and the total SIA
content is
$S_I = \sum_{n=1}^{i_{\mathrm{discrete}}} n\,c_n
       + \sum_{k=1}^{I_{\mathrm{bin}}} \mu_k^{(1)}$.
The number of moments per bin $P \in \{1,2,3\}$ is selected by the
\texttt{shape\_function} parameter:
\begin{itemize}
  \item \texttt{"constant"} (Eq.~\ref{eq:pc_ansatz}): track
        $\mu_k^{(0)}$ only ($P=1$).
  \item \texttt{"linear"} (Eq.~\ref{eq:shape_ansatz}): track
        $\bigl(\mu_k^{(0)}, \mu_k^{(1)}\bigr)$ ($P=2$).
  \item \texttt{"lognormal"} (Eq.~\ref{eq:lognormal_ansatz}): track
        $\bigl(\mu_k^{(0)}, \mu_k^{(1)}, \mu_k^{(2)}\bigr)$ ($P=3$).
\end{itemize}
The analogous bin moments
$\nu_k^{(p)} = \sum_{m \in \mathcal{B}^v_k} m^p\,\bar{c}_m$
are defined on the vacancy partition.

%-------------------------------------------------------------------
\subsubsection*{Bin-Integrated Rate Equations: ``Reconstruct -- Evaluate -- Project''}
\label{sec:bin_rate_eq}
%-------------------------------------------------------------------

Rather than re-deriving the bin-integrated rates symbolically (which
introduces explicit inter-bin flux terms and bin-averaged rate
constants), RadCluster\_1\_0 evolves the moments
$\{\mu_k^{(p)}\}$ via a three-step procedure executed at every
right-hand-side call (\texttt{bin\_moment\_rates.py}:
\texttt{reconstruct\_distribution}, \texttt{ode\_system},
\texttt{\_project\_to\_moments}):
\begin{enumerate}
  \item[\textbf{R.}] \textit{Reconstruct} a per-size distribution
        $c_n^{\mathrm{rec}}$ on the full integer grid
        $n = 1, \ldots, I$ by combining the discrete unknowns
        $\{c_n\}_{n \le i_{\mathrm{discrete}}}$ with the closure
        ansatz $c_n[\mu_k^{(0)},\mu_k^{(1)},\mu_k^{(2)}]$ for
        $n \in \mathcal{B}_k$ (Sec.~\ref{sec:closure} below). The
        same is done for the vacancy axis.
  \item[\textbf{E.}] \textit{Evaluate} the right-hand sides of the
        per-size master equations~\eqref{eq:ME_SIA},
        \eqref{eq:ME_vac}, \eqref{eq:ME_He}
        \emph{at every $n$ and every $m$}, using
        $c_n^{\mathrm{rec}}$ and $\bar{c}_m^{\mathrm{rec}}$. All
        reaction channels (production, coalescence, V--I annihilation,
        cavity absorption, thermal emission, fixed sinks, trap
        mutation, re-solution, He capture / emission) are evaluated
        \emph{without bin-averaging the rate constants}; the bin
        averaging enters only through the reconstructed concentrations.
  \item[\textbf{P.}] \textit{Project} the per-size rates onto the
        tracked moments:
        \begin{equation}\label{eq:projection}
          \frac{d\mu_k^{(p)}}{dt} \;=\;
          \sum_{n \in \mathcal{B}_k}
            n^p\,\Bigl(\frac{dc_n}{dt}\Bigr)^{\!\mathrm{full}},
          \qquad p = 0, 1, \ldots, P-1,
        \end{equation}
        and analogously for the vacancy moments. Sizes
        $n \le i_{\mathrm{discrete}}$ are copied verbatim into the
        discrete part of $d\mathbf{y}/dt$.
\end{enumerate}
This is the discrete analog of the Galerkin projection
\eqref{eq:galerkin} below: step \textbf{P} is the test-function
integration, with $n^p$ playing the role of $\phi_p(n)$ on the integer
grid $\mathcal{B}_k$. Because the per-size rate evaluation is
exact, all conservation properties of the master equations
(Eqs.~\ref{eq:ME_SIA}--\ref{eq:ME_He}) are inherited automatically: no
inter-bin flux $J_{k \to k+1}$ has to be defined or made
upwind-consistent, and no bin-averaged rate constants
$\hat{\mathcal{K}}_{k,m}^{(\mathrm{cav})}$ appear. The only
approximation is the closure used in step \textbf{R}.

\paragraph{Bin-integrated zeroth-moment equation.}
For documentation and analytic inspection, the projected
zeroth-moment equation reads, term-by-term,
\begin{equation}\label{eq:mu0_eq}
\begin{split}
    \frac{d\mu_k^{(0)}}{dt} = {} &
    \underbrace{
      \sum_{n \in \mathcal{B}_k} G_n
    }_{\text{cascade production}}
    %
    \;+\; \underbrace{
      \sum_{n \in \mathcal{B}_k}\!\Bigl[
        \tfrac{1}{2}\!\!\sum_{n'=1}^{n_{\max}^i}\!\!
          \mathcal{K}_{n',\,n-n'}^{ii}\,c_{n'}\,c_{n-n'}
        - c_n\!\!\sum_{n'=1}^{n_{\max}^i}\!
          \mathcal{K}_{n,n'}^{ii}\,c_{n'}
      \Bigr]
    }_{\text{SIA--SIA coalescence (gain + loss)}}
    \\[4pt]
    %--- thermal emission ---
    & \;+\; \underbrace{
      \sum_{n \in \mathcal{B}_k}\!\bigl[
        \varepsilon_i(n{+}1)\,c_{n+1}
        - \varepsilon_i(n)\,c_n\,\delta_{n \ge 2}
      \bigr]
    }_{\text{thermal SIA emission}}
    \\[4pt]
    %--- V-I gain ---
    & \;+\; \underbrace{
      \sum_{n \in \mathcal{B}_k}\sum_{m'=1}^{m_{\max}^v}
        \mathcal{K}_{m',n+m'}^{vi}\,\bar{c}_{m'}\,c_{n+m'}
      \;-\; \sum_{n \in \mathcal{B}_k}\sum_{m'=1}^{m_{\max}^v}
        \mathcal{K}_{m',n}^{vi}\,\bar{c}_{m'}\,c_n
    }_{\text{V--I annihilation (gain + loss); $n{=}1$ uses $\mathcal{K}_{iv}$}}
    \\[4pt]
    %--- cavity absorption ---
    & \;-\; \underbrace{
      \sum_{n \in \mathcal{B}_k}\sum_{m=1}^{V}
        \mathcal{K}_{n,m}^{(\mathrm{cav})}\,c_n\,\bar{c}_m
    }_{\text{SIA cluster--cavity absorption}}
    \;-\; \underbrace{
      \sum_{n \in \mathcal{B}_k} \mathcal{D}_i(n)\,c_n
    }_{\text{fixed-sink loss (mobile only)}}
    \\[4pt]
    & \;+\; \underbrace{
      \delta_{1 \in \mathcal{B}_k}\,
        \sum_{m,\ell}\Gamma_{\mathrm{TM}}(m,\ell)\,c_{m,\ell}
    }_{\text{trap mutation: only the bin containing $n{=}1$}}.
\end{split}
\end{equation}
The fixed-sink term retains its size dependence,
$\sum_n \mathcal{D}_i(n)\,c_n$, because
$\mathcal{D}_i(n) \propto \omega_n^{\mathrm{eff}}(n)$ varies with $n$
through the 1D-glide diffusivity (Sec.~\ref{sec:diffusion}).
The first-moment equation is the same projection weighted by $n$:
\begin{equation}\label{eq:mu1_eq}
    \frac{d\mu_k^{(1)}}{dt}
    \;=\; \sum_{n \in \mathcal{B}_k} n\,
        \Bigl(\frac{dc_n}{dt}\Bigr)^{\!\mathrm{full}},
\end{equation}
and the second moment, when tracked, follows the same pattern with
weight $n^2$.

\paragraph{Vacancy bin-integrated equations.}
The vacancy block obeys the analog of
Eq.~\eqref{eq:mu0_eq}--\eqref{eq:mu1_eq} with $n \to m$ and the
corresponding rate constants from~\eqref{eq:ME_vac}, including:
V--V coalescence, thermal vacancy emission
$\varepsilon_v(m,\ell)$ (with the He-pressure correction of
Sec.~\ref{sec:He_approx}), SIA-induced shrinkage for all mobile
$n = 1, \ldots, n_{\max}^i$, and the fixed-sink loss
$\sum_{m \le m_{\max}^v} \mathcal{D}_v\,\bar{c}_m$. He capture does
\emph{not} change the vacancy size class~$m$, so it appears in the
$dQ/dt$ equations only.

\paragraph{Helium equations.}
The two He-reduction modes of Sec.~\ref{sec:He_approx} are carried over
unchanged. In Case~2 (fission, decoupled scalar $\ell_{\mathrm{tot}}$)
a single scalar $Q_{\mathrm{tot}}$ is evolved alongside the moments;
in Case~1 (fusion, mean-field $\bar{\ell}(m)$) one $Q_k$ is attached
to each vacancy bin, so that
$\bar{\ell}_k = Q_k / \nu_k^{(0)}$ is constant within bin~$k$.

%-------------------------------------------------------------------
\subsection{Closure of bin-moment equations}
\label{sec:closure}
%-------------------------------------------------------------------

The reconstruction step \textbf{R} above requires a deterministic map
from the tracked moments $\bigl(\mu_k^{(0)}, \mu_k^{(1)}, \mu_k^{(2)}\bigr)$
back to per-size concentrations $c_n$ for $n \in \mathcal{B}_k$. Three
closures are implemented.

%-------------------------------------------------------------------
\paragraph{Piecewise-constant ($P{=}1$).}
\label{sec:pc_closure}
%-------------------------------------------------------------------
The simplest closure assumes the distribution is uniform within each
bin:
\begin{equation}\label{eq:pc_ansatz}
    c_n \;\approx\; \frac{\mu_k^{(0)}}{\Delta_k}
    \qquad \forall\; n \in \mathcal{B}_k.
\end{equation}
This is the classical \emph{logarithmic grouping}
\cite{Golubov2001,Ovcharenko2003,Golubov2007,Kohnert2017,Cui2020}. It
preserves $\mu_k^{(0)}$ exactly; the first-moment error per bin is
$O(\Delta_k/\hat{s}_k) = O(r-1)$. No first-moment ODE is solved in
this mode.

%-------------------------------------------------------------------
\paragraph{Linear dual basis ($P{=}2$, ``hat-function'').}
\label{sec:bin_galerkin}
%-------------------------------------------------------------------
With two tracked moments per bin, choose two basis functions
$\bigl(\phi_{k,0}(n), \phi_{k,1}(n)\bigr)$ such that
\begin{equation}\label{eq:shape_ansatz}
    c_n \;\approx\; \phi_{k,0}(n)\,\mu_k^{(0)}
                  + \phi_{k,1}(n)\,\mu_k^{(1)},
    \qquad n \in \mathcal{B}_k.
\end{equation}
The two basis functions are the linear \emph{dual basis} on the
discrete index set $\mathcal{B}_k$, i.e.\ the unique pair satisfying
\begin{equation}\label{eq:dual_basis_conditions}
  \sum_{n \in \mathcal{B}_k} \phi_{k,0}(n) = 1,
  \quad
  \sum_{n \in \mathcal{B}_k} \phi_{k,1}(n) = 0,
  \quad
  \sum_{n \in \mathcal{B}_k} n\,\phi_{k,0}(n) = 0,
  \quad
  \sum_{n \in \mathcal{B}_k} n\,\phi_{k,1}(n) = 1.
\end{equation}
Solving these constraints for the linear ansatz
$\phi_{k,p}(n) = \alpha_p + \beta_p\,n$ gives the closed form used in
the code:
\begin{equation}\label{eq:hat_basis_explicit}
    \phi_{k,0}(n) = \frac{S_2 - S_1\,n}{\mathcal{N}_k},
    \qquad
    \phi_{k,1}(n) = \frac{\Delta_k\,n - S_1}{\mathcal{N}_k},
\end{equation}
where the geometry integrals are
\begin{equation}\label{eq:hat_moments}
    S_1 \equiv \sum_{n \in \mathcal{B}_k} n,
    \qquad
    S_2 \equiv \sum_{n \in \mathcal{B}_k} n^2,
    \qquad
    \mathcal{N}_k \equiv \Delta_k\,S_2 - S_1^2 \;>\; 0,
\end{equation}
and are computed once at setup. Equations~\eqref{eq:hat_basis_explicit}
and \eqref{eq:hat_moments} are evaluated on the \emph{integer} grid
$\mathcal{B}_k = \{s_k, s_k+1, \ldots, s_{k+1}-1\}$; they reduce to
piecewise-linear FEM hat functions only in the continuum limit
$\Delta_k \to \infty$. Negative reconstructed values are clamped to
zero before use in step \textbf{E}; this clamp breaks exact moment
preservation only for strongly-skewed bins.

%-------------------------------------------------------------------
\paragraph{Log-normal closure ($P{=}3$).}
\label{sec:lognormal_closure}
%-------------------------------------------------------------------
Within each bin, assume the per-size concentration follows a
log-normal in $n$:
\begin{equation}\label{eq:lognormal_ansatz}
    c_n
    \;\approx\;
    \frac{a_k}{n}\,\exp\!\left[
      -\frac{(\ln n - m_k)^2}{2\,\sigma_k^2}
    \right],
    \qquad n \in \mathcal{B}_k,
\end{equation}
parameterised by $\bigl(a_k, m_k, \sigma_k^2\bigr)$. Given the three
tracked moments, define
$\bar{n}_k = \mu_k^{(1)}/\mu_k^{(0)}$ and
$\overline{n^2}_k = \mu_k^{(2)}/\mu_k^{(0)}$. The log-normal shape
parameters are obtained from the moment matching equations:
\begin{equation}\label{eq:lognormal_params}
    \sigma_k^2 \;=\; \ln\!\Biggl(
      \frac{\mu_k^{(2)}\,\mu_k^{(0)}}{\bigl(\mu_k^{(1)}\bigr)^{\!2}}
    \Biggr),
    \qquad
    m_k \;=\; \ln \bar{n}_k - \tfrac{1}{2}\,\sigma_k^2.
\end{equation}
The amplitude $a_k$ is fixed by requiring \emph{discrete} preservation
of the zeroth moment over $\mathcal{B}_k$:
\begin{equation}\label{eq:lognormal_norm}
    a_k \;=\; \frac{\mu_k^{(0)}}
        {\displaystyle
         \sum_{n \in \mathcal{B}_k}
           \frac{1}{n}
           \exp\!\Bigl[-\dfrac{(\ln n - m_k)^2}{2\,\sigma_k^2}\Bigr]}.
\end{equation}
Note: only $\mu_k^{(0)}$ is exactly preserved by this discrete
normalisation; $\mu_k^{(1)}$ and $\mu_k^{(2)}$ are recovered exactly
in the continuum limit but pick up
$O\!\bigl((\Delta_k / \hat{s}_k)^2\bigr)$ residuals on the integer
grid. When $\sigma_k^2 \le 0$ (monodisperse or non-physical
moments) the code falls back to the linear dual
basis~\eqref{eq:hat_basis_explicit}, and when $\Delta_k = 1$ to the
single-size assignment $c_{s_k} = \mu_k^{(0)}$.

%-------------------------------------------------------------------
\paragraph{Galerkin formulation (general).}
\label{sec:galerkin}
%-------------------------------------------------------------------
The three closures above are special cases of the general Galerkin
projection of the master equation~\eqref{eq:ME_SIA} onto basis
functions $\{\phi_p\}_{p=0}^{P-1}$ defined on $\mathcal{B}_k$:
\begin{equation}\label{eq:galerkin}
    \sum_{n \in \mathcal{B}_k}
      \phi_p(n)\,\frac{dc_n}{dt}
    \;=\;
    \sum_{n \in \mathcal{B}_k}
      \phi_p(n)\,f_n(\mathbf{c}),
    \qquad p = 0, 1, \ldots, P-1,
\end{equation}
with $f_n$ the right-hand side of~\eqref{eq:ME_SIA}. Choosing
$\phi_p(n) = n^p$ as the test functions yields the moment
projection~\eqref{eq:projection} used in the code; the trial
functions enter only through the reconstruction step \textbf{R}.

%-------------------------------------------------------------------
\subsubsection*{Accuracy and Conservation Properties}
\label{sec:moment_accuracy}
%-------------------------------------------------------------------

\paragraph{Mass conservation.}
Because the projection~\eqref{eq:projection} sums the
\emph{exact} per-size rates, the bin-moment system inherits the
Frenkel-pair and helium-balance identities of the full per-size system
without modification. In particular, the corrected swelling identity
(Sec.~\ref{sec:mass_conservation})
\begin{equation}\label{eq:swelling_identity_bin}
    S(t) \;=\; S_I(t)
              \;+\; \Delta J_i^{\mathrm{fix}}(t)
              \;-\; \Delta J_v^{\mathrm{fix}}(t),
\end{equation}
holds exactly at the level of the moment ODE system, with $S_I$ and
$S$ assembled as the discrete + first-moment sums introduced after
Eq.~\eqref{eq:bin_moments_def}. The associated diagnostic
\begin{equation}\label{eq:delta_FP_binned}
    \delta_{\mathrm{FP}}^{\mathrm{bin}}(t)
    \;=\;
    \frac{\bigl|S(t) - S_I(t)
         - \Delta J_i^{\mathrm{fix}}(t)
         + \Delta J_v^{\mathrm{fix}}(t)\bigr|}
         {S(t) + S_I(t)
         + \Delta J_i^{\mathrm{fix}}(t)
         + \Delta J_v^{\mathrm{fix}}(t)}
\end{equation}
is computed from the cumulative-loss integrals appended to the state
vector (\texttt{bin\_moment\_rates.py}: \texttt{i\_J\_SIA\_fixed},
\texttt{i\_J\_VAC\_fixed}, \texttt{i\_J\_He\_sink}), so it shares the
integrator's local-error control.

\paragraph{Truncation error.}
Closure error enters only through the reconstructed $c_n$ in step
\textbf{R}. For a smooth, slowly varying distribution
$f(n) = c_n$ inside a bin and any reaction kernel
$\mathcal{K}(n)$ that is smooth on the same scale:
\begin{equation}\label{eq:truncation_error}
    \sum_{n \in \mathcal{B}_k} \mathcal{K}(n)\,c_n
    \;-\;
    \sum_{n \in \mathcal{B}_k} \mathcal{K}(n)\,c_n^{\mathrm{rec}}
    \;=\;
    \begin{cases}
      O\!\bigl((r-1)^2\bigr), & P = 1\;\;\text{(piecewise-constant)}, \\
      O\!\bigl((r-1)^3\bigr), & P = 2\;\;\text{(linear dual basis)}, \\
      O\!\bigl((r-1)^4\bigr), & P = 3\;\;\text{(log-normal, smooth limit)}.
    \end{cases}
\end{equation}
For $r = 2$ this gives $\sim$25\%, $\sim$12\%, and $\sim$6\% per-bin
rate-constant error respectively; for $r = 1.5$, $\sim$6\%, $\sim$1.5\%,
$\sim$0.4\%. The lognormal closure is rate-limited by the integer-grid
residuals in $\mu_k^{(1)}, \mu_k^{(2)}$
(Eq.~\ref{eq:lognormal_norm}); in practice the linear dual basis with
$r \approx 1.5\text{--}2$ is the best accuracy / cost compromise and
is the default.

\paragraph{Recovery of the full system.}
Setting $I_{\mathrm{bin}} = 0, V_{\mathrm{bin}} = 0$ (i.e.\
$i_{\mathrm{discrete}} = I$, $v_{\mathrm{discrete}} = V$) recovers the
full per-size system: step~\textbf{R} becomes the identity, step
\textbf{E} is unchanged, and step \textbf{P} reduces to copying.
This is the regression check used to validate the bin-moment code
against \texttt{rate\_equations.py}.
